\section*{Chapter \ref{ch:b1:monopoly-pricing} --- Monopoly and Pricing}
% Numbers recomputed in code/bachelor-1/15-monopoly-pricing/exercises/solutions.py (tested).

\begin{solution}{exo:b1:monopoly-pricing:basic}
$MR = 100 - 2q = 20$: $q = 40$, $p = 60$, profit $40 \times 40 = 1600$. The competitive
quantity is 80, so the loss is $\tfrac12 \times 40 \times 40 = 800$.
\end{solution}

\begin{solution}{exo:b1:monopoly-pricing:lerner-ex}
$L = 1/4 = 0.25$ and $\mu = 1/(1 - 0.25) = 1.333$, or $\varepsilon/(1 + \varepsilon) = -4/-3$.
\end{solution}

\begin{solution}{exo:b1:monopoly-pricing:read}
The gap is $32.04 - 2.98 = 29.06$ points. Other reasons: more loans that are not
repaid, so a larger risk premium; higher and more uncertain inflation;
reserve requirements and taxes on banking that raise the cost of lending;
higher operating costs per loan; and differences in how the two rates are
defined in each country's statistics.
\end{solution}

\begin{solution}{exo:b1:monopoly-pricing:predict}
Tourists: the less elastic group pays more
(\cref{prop:b1:monopoly-pricing:third}).
\end{solution}

\begin{solution}{exo:b1:monopoly-pricing:third-ex}
Market 1: $100 - 2q_1 = 20$, $q_1 = 40$, $p_1 = 60$. Market 2: $60 - 2q_2 = 20$, $q_2 = 20$,
$p_2 = 40$. Single price (both served, $p \le 60$): $Q = 160 - 2p$, $MR = 80 - Q = 20$, so $Q = 60$
and $p = 50$, with $q_1 = 50$ and $q_2 = 10$. (Serving market 1 alone at 60 earns 1,600,
less than $30 \times 60 = 1800$.)
\end{solution}

\begin{solution}{exo:b1:monopoly-pricing:welfare}
Discrimination: \omterm{def:b1:consumer-welfare:consumer-surplus}{consumer surplus} $\tfrac12 \times 40 \times 40 + \tfrac12 \times 20 \times 20 = 1000$, profit
$40 \times 40 + 20 \times 20 = 2000$, total 3,000. Single price: $\tfrac12 \times 50 \times 50 + \tfrac12 \times 10 \times 10 = 1300$
and $30 \times 60 = 1800$, total 3,100. The most is $\tfrac12 \times 80 \times 80 + \tfrac12 \times 40 \times 40 = 4000$, so the
\omterm{def:b1:partial-equilibrium-welfare:dwl}{deadweight loss} is 1,000 under discrimination and 900 under the single price.
\end{solution}

\begin{solution}{exo:b1:monopoly-pricing:tpt}
Unit price 2, fee $\tfrac12 \times 8 \times 8 = 32$ per buyer.
\end{solution}

\begin{solution}{exo:b1:monopoly-pricing:bundle}
Separately, a price of 8 sells each good once (16) and a price of 3 twice
(12): 16. A bundle at 11 sells to both: 22. With values $(8, 8)$ and $(3, 3)$
the bundle is worth 16 and 6: at 16 it earns 16, at 6 it earns 12, no more
than selling separately. \omterm{def:b1:monopoly-pricing:bundling}{Bundling} pays when values are negatively
correlated.
\end{solution}

\begin{solution}{exo:b1:monopoly-pricing:cap}
The cap falls by about $1\%$ a year in money: $0.99^5 = 0.951$, about $4.9\%$ lower
after five years, and $1 - (0.99/1.02)^5 = 13.9\%$ lower in real terms (about $3\%$ a year). Since the company
keeps any cost it saves until the next review, it cuts costs faster than
under \omterm{def:b1:monopoly-pricing:regulation}{average-cost pricing}, where savings passed to buyers at once.
\end{solution}

\begin{solution}{exo:b1:monopoly-pricing:derive}
$R(q) = p(q)q$, so $MR = p + qp'(q) = p(1 + \frac qp\frac{\dd p}{\dd q}) = p(1 + 1/\varepsilon)$. Setting it equal
to $MC$ gives $(p - MC)/p = -1/\varepsilon$. If $-1 < \varepsilon < 0$, $1 + 1/\varepsilon < 0$, so $MR < 0 < MC$:
lowering output raises revenue and lowers cost.
\end{solution}

\begin{solution}{exo:b1:monopoly-pricing:ramsey}
$(p - 20)(100 - p) = 1200$ gives $p^2 - 120p + 3200 = 0$, so $p = 40$ (or 80), $q = 60$ and
the loss is $\tfrac12 \times 20 \times 20 = 200$. Any lower price loses money, and the loss
grows with the price, so the lowest break-even price is the best. With
several goods the second-best prices (Ramsey prices) set each good's \omterm{def:b1:monopoly-pricing:lerner}{Lerner
index} proportional to the inverse of its elasticity: the highest \omterm{def:b1:monopoly-pricing:markup}{markup} on
the least elastic good.
\end{solution}

\begin{solution}{exo:b1:monopoly-pricing:refute}
If the rise reflects fixed costs, profit rates net of all costs should not
rise, entry should be as easy as before, and \omterm{def:b1:monopoly-pricing:markup}{markups} measured with overhead
treated as a variable cost should not rise. Rising profit rates, falling
entry, and a rise that survives the change in cost measure all point to
market power; their absence would refute it.
\end{solution}

\begin{solution}{pb:b1:monopoly-pricing:ferry}
\textbf{1.} $Q = 240 - 5p$ for $p \le 40$ (only tourists above 40).

\textbf{2.} $p = 48 - Q/5$, $MR = 48 - 2Q/5 = 10$: $Q = 95$ and $p = 29$.

\textbf{3.} Residents $160 - 116 = 44$, tourists $80 - 29 = 51$ thousand trips.

\textbf{4.} $19 \times 95 = 1805$. Tourists alone: $80 - 2q = 10$, $q = 35$, $p = 45$, profit
$35 \times 35 = 1225 < 1805$.

\textbf{5.} $L = 19/29 = 0.655$; $\varepsilon = -5 \times 29/95 = -1.526$ and $1/1.526 = 0.655$.

\textbf{6.} Residents $\tfrac12 \times 44 \times 11 = 242$; tourists $\tfrac12 \times 51 \times 51 = 1300.5$.

\textbf{7.} At $p = 10$ residents make 120 and tourists 70. Residents
$\tfrac12 \times 19 \times 76 = 722$, tourists $\tfrac12 \times 19 \times 19 = 180.5$: 902.5.

\textbf{8.} $p = 40 - q/4$, $MR = 40 - q/2 = 10$: $q = 60$, $p = 25$.

\textbf{9.} $MR = 80 - 2q = 10$: $q = 35$, $p = 45$.

\textbf{10.} Both are 10. $\varepsilon_R = -4 \times 25/60 = -1.667$ and $\varepsilon_T = -45/35 = -1.286$: the more
elastic residents pay less.

\textbf{11.} $15 \times 60 + 35 \times 35 = 900 + 1225 = 2125$.

\textbf{12.} Residents $\tfrac12 \times 60 \times 15 = 450$; tourists $\tfrac12 \times 35 \times 35 = 612.5$.

\textbf{13.} Residents $\tfrac12 \times 15 \times 60 = 450$, tourists $\tfrac12 \times 35 \times 35 = 612.5$: 1,062.5, more than
902.5.

\textbf{14.} With linear demands each group gets half its competitive
quantity, $60 + 35 = 95$, as under the single price. But tourists willing to pay
between 29 and 45 lose trips, and residents willing to pay between 25 and
29 gain them: trips move to buyers who value them less.

\textbf{15.} $\tfrac12 \times 30 \times 30 = 450$ dalens a year.

\textbf{16.} Fee 450 dalens a year and 10 a trip; each resident makes 30 trips,
120 thousand in all.

\textbf{17.} $4000 \times 450 = 1.8$ million dalens, 1,800; with tourists, $1800 + 1225 = 3025$.

\textbf{18.} Residents now make every trip worth its cost: the loss is the
tourists', 612.5.

\textbf{19.} Tourists differ and most cross once or twice; the operator
cannot see each one's demand, and a pass sold to a few tourists could be
passed on to others.

\textbf{20.} $(902.5 - 612.5)/902.5 = 32.13\%$.

\textbf{21.} $C(Q) = 1500 + 10Q$, and
$C(Q_1) + C(Q_2) = 3000 + 10(Q_1 + Q_2) > C(Q_1 + Q_2)$: one ferry is always cheaper.

\textbf{22.} 190 thousand trips; revenue pays only the marginal costs, and the
ferry loses the 1,500.

\textbf{23.} $(p - 10)(240 - 5p) = 1500$ gives $p^2 - 58p + 780 = 0$, so
$p = (58 - \sqrt{244})/2 = 21.19$ and $Q = 134.05$.

\textbf{24.} Residents make $160 - 84.76 = 75.24$: $\tfrac12 \times 11.19 \times 44.76 = 250.4$; tourists
$\tfrac12 \times 11.19 \times 11.19 = 62.6$: 313.0.

\textbf{25.} $21.19 \times (1 + 0.02 - 0.03) = 20.98$.

\textbf{26.} $Q = 240 - 5 \times 20.98 = 135.1$ and profit $(20.98 - 9) \times 135.1 - 1500 = 118.5$ with
these rounded values (118.3 unrounded). Under \omterm{def:b1:monopoly-pricing:regulation}{average-cost pricing} the price
would be reset to the new average cost and the profit would be zero.

\textbf{27.} The residents' season pass removes $32.1\%$ of the single-price
\omterm{def:b1:partial-equilibrium-welfare:dwl}{deadweight loss}: 290 of 902.5.
\end{solution}
